\question[4]
Erläutere, was es bedeutet, dass eine Datenbank in der 2. Normalform vorliegt.
Gib ein Beispiel an für eine Datenbank, die nicht in der 2. Normalform vorliegt,
und erkläre, wie man sie in die 2. Normalform überführen kann.


\begin{solutionbox}{12cm}
Eine Datenbank befindet sich in der 2. Normalform, wenn sie in der 1. Normalform vorliegt
und jedes Nicht-Schlüsselattribut voll vom gesamten Schlüssel abhängt.

\vspace{0.2cm}

\textbf{Beispiel für keine 2. Normalform:}

\begin{center}
\begin{tabular}{|l|l|l|}
\hline
KursID & SchülerID & SchülerName \\
\hline
10 & 1 & Anna Meier \\
10 & 2 & Max Schulz \\
11 & 1 & Anna Meier \\
\hline
\end{tabular}
\end{center}

Hier ist der Schlüssel zusammengesetzt aus \textit{KursID} und \textit{SchülerID}.
Das Attribut \textit{SchülerName} hängt aber nur von \textit{SchülerID} ab und damit nicht vom ganzen Schlüssel.

\vspace{0.2cm}

\textbf{Überführung in die 2. Normalform:}

\begin{center}
\begin{tabular}{|l|l|}
\hline
SchülerID & SchülerName \\
\hline
1 & Anna Meier \\
2 & Max Schulz \\
\hline
\end{tabular}
\end{center}

\begin{center}
\begin{tabular}{|l|l|}
\hline
KursID & SchülerID \\
\hline
10 & 1 \\
10 & 2 \\
11 & 1 \\
\hline
\end{tabular}
\end{center}

Nun sind alle Nicht-Schlüsselattribute voll vom gesamten Schlüssel abhängig.
Die Tabellen liegen damit in der 2. Normalform vor.

\end{solutionbox}
